\documentclass[11pt]{article}
\usepackage[utf8]{inputenc}
\usepackage[english]{babel}
\usepackage{geometry}
\usepackage{amsmath}
\usepackage{booktabs}
\usepackage{xcolor}
\usepackage{mdframed}

\geometry{a4paper, margin=1in}

\title{Technical Specifications: HackRF Pro}
\author{Radiofrequency Engineering}
\date{\today}

\newmdenv[
  backgroundcolor=red!8,
  linecolor=red!60,
  linewidth=1pt,
  roundcorner=4pt,
  innertopmargin=6pt,
  innerbottommargin=6pt
]{warningbox}

\newmdenv[
  backgroundcolor=yellow!10,
  linecolor=yellow!60!black,
  linewidth=1pt,
  roundcorner=4pt,
  innertopmargin=6pt,
  innerbottommargin=6pt
]{pendingbox}

\begin{document}

\maketitle
\tableofcontents
\newpage

% =========================================================
\section{Gains}
% =========================================================

Controlled separately through \texttt{libhackrf}:

\begin{itemize}
    \item \textbf{LNA (IF RX):} \texttt{hackrf\_set\_lna\_gain()} - 0 to 40 dB, 8 dB steps.
    \item \textbf{VGA (baseband RX):} \texttt{hackrf\_set\_vga\_gain()} - 0 to 62 dB, 2 dB steps.
    \item \textbf{TX VGA:} \texttt{hackrf\_set\_txvga\_gain()} - 0 to 47 dB, 1 dB steps.
    \item \textbf{AMP (RF front end):} \texttt{hackrf\_set\_amp\_enable()} - 0 (off) or 1 (on, fixed $\approx$+14 dB). Applies to RX and TX.
\end{itemize}

\textbf{Note:} Adjust gain while avoiding ADC saturation; validate with the noise floor and absence of clipping.

% =========================================================
\section{Hardware Filters}
% =========================================================

\begin{itemize}
    \item \textbf{Analog filter (Baseband LPF):} \texttt{hackrf\_set\_baseband\_filter\_bandwidth()}\\
    Selectable thresholds: 1.75, 2.5, 3.5, 5, 5.5, 14, 20, 24, 28 MHz.\\
    Transceiver: MAX2831 (Pro) - compatible with existing software through \textit{legacy mode} in \texttt{libhackrf}.

    \item \textbf{Digital filtering and decimation (iCE40 FPGA):}\\
    Configured implicitly according to the sample rate. Handles hardware decimation, removes the DC spike (previously handled in software), and enables 16-bit samples at low rates with a typical ENOB of 9--11 bits.
\end{itemize}

% =========================================================
\section{Frequency Ranges}
% =========================================================

\begin{itemize}
    \item \textbf{Nominal range:} 100 kHz to 6 GHz.
    \item \textbf{Tunable range (physical limits):} 0 Hz to 7.1 GHz.
    \begin{itemize}
        \item Below 100 kHz: analog sensitivity drops drastically.
        \item Above 6 GHz: severe attenuation and unreliable calibration.
    \end{itemize}
\end{itemize}

% =========================================================
\section{Sample Rate and Resolution}
% =========================================================

\begin{table}[h]
    \centering
    \small
    \begin{tabular}{@{}p{0.20\textwidth}p{0.15\textwidth}p{0.22\textwidth}p{0.34\textwidth}@{}}
        \toprule
        \textbf{Mode} & \textbf{Resolution} & \textbf{Max. Sample Rate} & \textbf{Typical Use} \\
        \midrule
        Standard       & 8-bit I/Q & 20 MS/s & General compatibility, up to 20 MHz of visible BW \\
        Extended precision & 16-bit & Low (FPGA decimation) & Narrowband signals, higher dynamic range (ENOB 9--11) \\
        Medium precision & 4-bit & 40 MS/s & Wideband sweep; reduced resolution to fit over USB \\
        \bottomrule
    \end{tabular}
\end{table}

% =========================================================
\section{IQ \& DC Corrections}
% =========================================================

\begin{itemize}
    \item \textbf{DC spike:} Removed in hardware by the iCE40 FPGA before sending data over USB. It no longer requires software filters (GNU Radio, etc.).
    \item \textbf{I/Q imbalance:}
    \begin{pendingbox}
    Not documented in \texttt{libhackrf} or in the official firmware to date. Verify in future API or FPGA versions.
    \end{pendingbox}
\end{itemize}

% =========================================================
\section{Connectors and Inputs}
% =========================================================

\begin{table}[h]
    \centering
    \small
    \begin{tabular}{p{0.18\textwidth}p{0.26\textwidth}p{0.44\textwidth}}
        \toprule
        \textbf{Connector} & \textbf{Function} & \textbf{Specification} \\
        \midrule
        USB-C      & Power and data & USB 2.0 High Speed. 5 V / 500 mA min. \\
        RF (SMA)   & Main TX/RX & 50 $\Omega$. Half-duplex. Range: 100 kHz -- 7.1 GHz. \\
        P1 (SMA)   & CLK IN (default) & 10 MHz, 3.3 V square wave. Do not exceed 3.3 V or go below 0 V. \\
        P2 (SMA)   & CLK OUT (default) & 10 MHz, 3.3 V. Daisy chain between devices. \\
        P20        & GPIO / ADC / RTC & 22 pins, 3.3 V. GPIO, ADC, and auxiliary power access. \\
        P22        & I2S / SPI / I2C / UART & 26 pins. Alternate CLK IN (enable with \texttt{hackrf\_clock -c p22}). \\
        P28        & SD / Trigger & 22 pins. SDIO, trigger IN/OUT for synchronization. \\
        \bottomrule
    \end{tabular}
\end{table}

\begin{itemize}
    \item \textbf{Integrated TCXO:} 0.5 ppm. Eliminates frequency drift without an external reference.
    \item \textbf{External clock:} Connect 10 MHz to P1. The hardware switches automatically when TX/RX starts. Verify detection with \texttt{hackrf\_clock -i}.
    \item \textbf{Clock out:} Enable with \texttt{hackrf\_clock -o 1}. Allows daisy chaining multiple HackRF Pro units.
    \item \textbf{RF Bias-T:} 3.3 V, 50 mA max. Software-enabled to power external LNAs.
    \item \textbf{GPIO:} All headers operate at 3.3 V CMOS. P22 includes I2C and UART for external peripheral control.
    \item \textbf{Note:} The audio jack and MicroSD slot are part of the \textbf{PortaPack H2} module, not the base HackRF Pro.
\end{itemize}

% =========================================================
\section{Safe Operating Limits}
% =========================================================

\begin{warningbox}
\textbf{Critical warnings - read before connecting external signals}
\end{warningbox}

\vspace{4pt}
\begin{itemize}
    \item \textbf{Maximum RX power:} Do not exceed \textbf{-5 dBm} to avoid distortion. Irreversible physical damage starts at \textbf{+10 dBm} (even with AMP disabled, a software error can enable it).
    \item \textbf{TX$\to$RX loopback:} Use sufficient attenuation. The TX level can reach 10--15 dBm; without attenuation the front end is damaged.
    \item \textbf{Unloaded RF port:} Always terminate with an antenna or \textbf{50 $\Omega$} during transmission.
    \item \textbf{External clock P1:} 10 MHz signal strictly between 0 V and 3.3 V. Do not use other frequencies except with firmware modification.
    \item \textbf{GPIO:} Do not exceed 3.3 V on any pin of headers P20/P22/P28.
    \item \textbf{ESD:} Handle with antistatic protection. Do not allow metal contact with the powered PCB.
\end{itemize}

% =========================================================
\section{Clock and Synchronization}
% =========================================================

\begin{itemize}
    \item \textbf{Internal clock:} 0.5 ppm TCXO, active by default. Requires no configuration.
    \item \textbf{External clock:} 10 MHz signal on P1 (SMA). The device switches automatically when a TX/RX operation starts. Verify with \texttt{hackrf\_clock -i}.
    \item \textbf{Clock out:} 10 MHz signal available on P2. Enable with \texttt{hackrf\_clock -o 1}. Useful for daisy chaining units or synchronizing with laboratory instruments.
    \item \textbf{Alternate signals on P1/P2:} Using \texttt{hackrf\_clock -1} or \texttt{hackrf\_clock -2}, different internal signals can be connected instead of the defaults (trigger, FPGA clocks, etc.).
    \item \textbf{Alternate CLK IN:} P22 pin 2 exposes a second independent CLKIN input. Enable with \texttt{hackrf\_clock -c p22}. Do not connect P1 and P22 pin 2 simultaneously as inputs.
    \item \textbf{Trigger IN/OUT:} Available on P28 pins 15 and 16. Allows capture-start synchronization between devices.
    \item \textbf{Multi-HackRF:} Connect CLKOUT from one unit to CLKIN of the next. Enable CLKOUT with \texttt{hackrf\_clock -o 1} before starting operation.
\end{itemize}

% =========================================================
\section{Indicator LEDs}
% =========================================================

When connecting the HackRF Pro to USB, \textbf{four LEDs should turn on}: MCU, FPGA, RF, and USB. If any of them does not turn on, it indicates a fault in the internal power chain or in the firmware.

\begin{table}[h]
    \centering
    \begin{tabular}{p{0.12\textwidth}p{0.36\textwidth}p{0.40\textwidth}}
        \toprule
        \textbf{LED} & \textbf{Function} & \textbf{Normal state} \\
        \midrule
        MCU  & Primary internal power and active firmware & On when USB is connected \\
        FPGA & Active iCE40 FPGA power                   & On when USB is connected \\
        RF   & Active RF-chain power                     & On when USB is connected \\
        USB  & USB communication established with host   & On when USB is connected \\
        RX   & Reception operation in progress           & On only during active RX \\
        TX   & Transmission operation in progress        & On only during active TX \\
        \bottomrule
    \end{tabular}
\end{table}

\textbf{Note:} LED colors have no meaning - they differ from one another to make visual distinction easier, not to indicate states.

% =========================================================
\section{Expansion Pinout (P20 / P22 / P28)}
% =========================================================

The HackRF Pro exposes three internal expansion headers. All operate at \textbf{3.3 V CMOS}.

\subsection*{P20 - GPIO / ADC / Power (22 pins)}

\begin{table}[h]
    \centering
    \small
    \begin{tabular}{clc|clc}
        \toprule
        \textbf{Pin} & \textbf{Signal} & & \textbf{Pin} & \textbf{Signal} & \\
        \midrule
        1  & VBAT       & & 2  & PB\_5        & \\
        3  & 3V3AUX     & & 4  & WAKEUP       & \\
        5  & GPIO3\_8   & & 6  & GPIO3\_9     & \\
        7  & GPIO3\_10  & & 8  & GPIO3\_11    & \\
        9  & GPIO3\_12  & & 10 & GPIO3\_13    & \\
        11 & GPIO3\_14  & & 12 & GPIO3\_15    & \\
        13 & GND        & & 14 & ADC0\_6      & \\
        15 & GND        & & 16 & ADC0\_2      & \\
        17 & VBUSCTRL   & & 18 & ADC0\_5      & \\
        19 & GND        & & 20 & ADC0\_0      & \\
        21 & VBUS       & & 22 & VIN          & \\
        \bottomrule
    \end{tabular}
\end{table}

\subsection*{P22 - I2S / SPI / I2C / UART (26 pins, selection)}

Includes: alternate CLKIN (pin 2), I2C1 (pins 5--6), UART (pins 19--20), I2S TX/RX, and FPGA auxiliary signals. See the full documentation at \texttt{hackrf.readthedocs.io}.

\subsection*{P28 - SD / Trigger / GPIO (22 pins)}

Includes: full SDIO (pins 3--11), Trigger OUT (pin 15), and Trigger IN (pin 16) for hardware synchronization between devices. See the full documentation at \texttt{hackrf.readthedocs.io}.

\textbf{Usage note:} Do not exceed 3.3 V on any pin. Confirm maximum current per pin in the LPC43xx datasheet before connecting external loads.

\end{document}
